\section{Related Work}\label{sec:related}

\para{Linear representations: single direction vs.\ subspace.}
A large body of work shows that high-level concepts are encoded as linear directions in LLM activation space. \citet{park2023linear} formalize the linear representation hypothesis, \citet{tigges2023linear} localize sentiment to a single direction, and \citet{arditi2024refusal} show that refusal is mediated by one direction. \citet{marks2023geometry} establish that difference-in-means directions are both robust and causal, the estimator we adopt throughout. A competing line argues that many concepts occupy a multi-dimensional \emph{subspace} rather than a single axis: \citet{wollschlager2025geometry} describe concept cones, \citet{zhao2024gcs} fit Gaussian concept subspaces, \citet{shafran2026regions} characterize concepts as regions amenable to multi-feature analysis, and \citet{shah2025harmfulness} probe subconcepts to recover a low-rank harmfulness subspace, a template we adapt to build the drama subspace. \citet{subramani2022steering} steer generation along extracted directions. Unlike this work, which adjudicates single-axis versus subspace on factual or behavioral concepts, we test the geometry of a fresh, evaluative, reader-dependent construct, and we do so \emph{separately} for the construct itself and for the variation across readers.

\para{Structured and relational probes.}
Beyond single directions, probes can recover structured and relational geometry. \citet{hewitt2019structural} extract syntax trees, \citet{nanda2023emergent} find emergent world models, and \citet{marks2024sfc} attribute behavior through sparse feature circuits. \citet{dawes2026hprobes} probe conceptual hierarchy, \citet{kozlowski2026semantic} show that evaluative axes live on a shared low-dimensional subspace with cross-axis spillover, and \citet{lee2026tpr} factorize representations along a reader $\times$ content structure. These results directly motivate our central test: whether reader variation factorizes as a reader $\times$ drama interaction (H2). Building on this framing, we operationalize and test that factorization explicitly, and find it absent.

\para{Affect geometry in LLMs.}
Closest in subject matter, recent work maps the geometry of affect inside LLMs. \citet{reichman2025emotions} find a low-dimensional, steerable emotional manifold and introduce a synthetic-to-real transfer protocol that we reuse to validate our controlled corpus against naturalistic data. \citet{wang2025emotion} localize emotion circuits, and \citet{zhao2025hierarchical} describe a hierarchical, persona-biased emotional representation. This evidence supports the low-dimensional premise underlying H1. We extend it with a reader-conditioning test that asks not just how affect is represented, but how the representation of what is \emph{dramatic} varies across the people perceiving it.

\para{Perspective and reader axes.}
A reader's viewpoint can itself be a low-dimensional direction. \citet{kim2025political} recover a political-perspective axis, and \citet{qin2026granularity} find a role-granularity axis governing how models adopt perspectives. Unlike these studies, which establish that a perspective axis exists, we ask whether such a reader axis \emph{hosts} a drama subspace, that is, whether different readers read drama along different directions or merely set different thresholds on a shared one.

\para{Subjectivity and per-reader modeling.}
Annotation-disagreement research treats reader variation as signal rather than noise. \citet{song2025survey} survey learning with disagreement; \citet{orlikowski2025beyond} show that reader effects are largely idiosyncratic rather than demographic; and \citet{kanclerz2022personalized,kanclerz2023pals,bielaniewicz2023chuckles} build personalized models of subjective phenomena including humor. \citet{mehrotra2026multi}, \citet{zhou2025subjectivity}, \citet{mieleszczenko2023capturing}, and \citet{alexeeva2023annotating} further model per-annotator labels and contested items. This work establishes that reader identity carries large, learnable variance that is not reducible to low-dimensional demographics, which we confirm in the strong per-reader threshold effect. Adding to it, we show in representation space that this variance is also not a reader-specific drama \emph{direction}: readers differ in where they place a threshold, not in which way they read.

\para{Narrative drama and reader response.}
Finally, computational narrative work grounds our construct. \citet{tian2024narratives} and \citet{christ2024emotional} model arousal and suspense trajectories, \citet{subbiah2025spinning} study narrative ambiguity, and \citet{dai2026cedar} document cross-cultural divergence in affective response. We operationalize drama as perceived emotional intensity in this tradition, drawing labels from \goemotions \citep{demszky2020goemotions} and inspecting \emobank \citep{buechel2017emobank} as auxiliary aggregate evidence. Unlike this work, which studies drama and reader response behaviorally, we connect it to the internal geometry of an LLM.
